% TOP_PROBLEM: {"id": 30004279, "title": "Equivalent Bicommutant Categories from Nonisomorphic Conformal Nets", "category_id": 16, "category": {"id": 16, "name": "physics", "display_name": "Mathematical Physics", "description": "Problems at the intersection of mathematics and physics.", "slug": "physics", "order_index": 16, "created_at": "2026-07-31T15:26:25.671Z"}, "status": "open", "statusQualification": "The size and parameterization of large families and the precise associated category must be the report\\textquotesingle s definitions. The question concerns equivalence of that category, not merely equality of numerical invariants or representation-ring dimensions."}
% FIRST_POSTED: 2026-10-11
% ENTRY_KIND: statement_only
% SUBMISSION: {"kind": "statement_only", "category": "top_problems", "problem_id": "30004279", "proposal_type": "established", "source": "Subfactors and Applications (2019), Oberwolfach Reports, DOI 10.4171/owr/2019/49, pp. 3053--3140. Primary report accessed 2026-10-11; relevant PDF pages 23, 24. https://doi.org/10.4171/owr/2019/49; Henriques--Penneys, representations of fusion categories and commutants (catalogue-listed context; not independently consulted in this preparation). https://arxiv.org/abs/2004.08271; Henriques--Penneys--Tener, finite-depth objects in bicommutant categories (catalogue-listed context; not independently consulted in this preparation). https://arxiv.org/abs/2307.13822", "status_evidence": "Canonical source status open retained. Supplied literature triage dated 2026-08-22: Morita-equivalent unitary fusion categories are known to yield equivalent bicommutant categories, and finite-depth objects admit later classifications, but the stated existence of many non-isomorphic conformal nets with equivalent bicommutant tensor categories is not established by the cited literature. Primary-report access checked 2026-10-11; this historical report and inherited triage do not certify current resolution of every component.", "status_checked_on": "2026-10-11", "duplicate_check": "Compared canonical ID, title and normalized target with the complete local UnsolvedMath registry; normalized exact-text matches: none; title matches: none. Related or overlapping canonical records: none identified in this batch. Primary source context and the existing OWR statement audit used for scope and duplicate review; an exhaustive mathematical-equivalence/literature search was not performed.", "attribution": "Original proposer not named in the supplied canonical record; attribution remains to the identified Oberwolfach contribution/report, not its editors as invented proposers. Explanatory statement preparation and source/semantic audit assisted by OpenAI GPT-6 Codex; no human mathematical authorship or independent proof verification is claimed."}
% SOURCE_STATEMENT: "Do there exist large families of pairwise nonisomorphic conformal nets whose associated bicommutant categories are equivalent as tensor categories?"
% Preparation and final self-audit: CONTRIBUTING.md prompts C and D; source/audit artifacts retained in the local physics_statements_2026_10_11 batch.
% Imported from: physics_statements.tex; UnsolvedMath ID 30004279
% !TeX program = lualatex
\documentclass[11pt]{article}
\usepackage[margin=25mm]{geometry}
\usepackage{fontspec}
\setmainfont{Latin Modern Roman}
\usepackage{amsmath,amssymb,mathtools}
\usepackage{unicode-math}
\setmathfont{Latin Modern Math}
\usepackage{xurl,hyperref}
\hypersetup{colorlinks=true,urlcolor=blue}
\setcounter{secnumdepth}{0}
\setlength{\parindent}{0pt}
\setlength{\parskip}{5pt}
\setlength{\emergencystretch}{3em}
\newcommand{\sref}[2]{\hyperlink{source:#1:#2}{[S#2]}}
\newcommand{\cataloguescope}[1]{\paragraph{Recorded target.}#1}
\newcommand{\R}{\mathbb R}
\newcommand{\C}{\mathbb C}
\newcommand{\N}{\mathbb N}
\newcommand{\Z}{\mathbb Z}
\newcommand{\Q}{\mathbb Q}
\newcommand{\F}{\mathbb F}
\newcommand{\eps}{\varepsilon}
\begin{document}
% UnsolvedMath ID 30004279; source catalogue code OWR-17292-002
\setcounter{section}{30004278}
\section{Equivalent Bicommutant Categories from Nonisomorphic Conformal Nets}\label{30004279}
{\small UnsolvedMath ID 30004279 · Mathematical Physics}
\subsection{Definitions and mathematical statement}

\paragraph{Shared notation.}$\mathbb N=\{1,2,\ldots\}$; the standard number fields have their usual meanings. All additional conventions are stated locally.


\paragraph{Definitions and conventions.}
A conformal net assigns local von Neumann algebras to intervals, with locality and conformal covariance. An associated bicommutant category is the tensor category constructed from the net by the source\textquotesingle s operator-algebraic analogue of taking a double commutant. Tensor equivalence is equivalence respecting products and unit objects. Nonisomorphic nets can differ as local theories even when a derived category fails to distinguish them.

\paragraph{Target retained from the canonical source catalogue.}
\paragraph{Catalogue wording retained for traceability.}
Do there exist large families of pairwise nonisomorphic conformal nets whose associated bicommutant categories are equivalent as tensor categories?

\paragraph{Source scope and qualification.}
The size and parameterization of large families and the precise associated category must be the report\textquotesingle s definitions. The question concerns equivalence of that category, not merely equality of numerical invariants or representation-ring dimensions.
\subsection{Short English statement}
Can large families of distinct conformal nets share equivalent associated bicommutant tensor categories?
\subsection{Sources}
\begin{itemize}
\item[\hypertarget{source:30004279:1}{S1}] ULAM.AI and the UnsolvedMath Contributors, UnsolvedMath, record 30004279 (OWR-17292-002). Catalogue curation under CC BY 4.0; definitions and English exposition added locally.
\url{https://huggingface.co/datasets/ulamai/UnsolvedMath}.
\item[\hypertarget{source:30004279:2}{S2}] Subfactors and Applications (2019), Oberwolfach Reports, DOI 10.4171/owr/2019/49, pp. 3053--3140. Primary report accessed 2026-10-11; relevant PDF pages 23, 24.
\url{https://doi.org/10.4171/owr/2019/49}.
\item[\hypertarget{source:30004279:3}{S3}] Henriques--Penneys, representations of fusion categories and commutants (catalogue-listed context; not independently consulted in this preparation).
\url{https://arxiv.org/abs/2004.08271}.
\item[\hypertarget{source:30004279:4}{S4}] Henriques--Penneys--Tener, finite-depth objects in bicommutant categories (catalogue-listed context; not independently consulted in this preparation).
\url{https://arxiv.org/abs/2307.13822}.
\end{itemize}
\label{end:30004279}
\end{document}
